\scanpage{135}
\uline{de dégénérescence}, \uline{alors} (219) \uline{est \unsure{un|une}}
\deleted{\unsure{équiv.}}\uline{isomorphisme}, de nature ess.t triviale.
On est donc réduit à regarder le
cas d'une application \uline{injective} $\alpha$ (« opération
face »), et le \inserted{seul} cas qui \deleted{\unsure{pose}} pose problème
est bien sûr celui où $\alpha(0)\ne0$. Donc
\[
\tag{221}
\left\{\begin{aligned}
d_*&=\{d_0\unsure{\geq|>}d_1\cdots>d_r\},\quad\text{et si}\\
d'_*&=\{d'_0>d'_1\unsure{\geq|>}\cdots>d'_{r'}\}\text{ est }
\deleted{\unsure{une sous-suite}}\inserted{\text{contenu dans }d_*},\\
&\quad d'_0\ne d_0\text{ i.e. }d'_0<d_0,\quad\text{il faut}
\end{aligned}\right.
\]
définir l'hom
\[
(222)
\]
\diagram{\begin{tikzcd}[ampersand replacement=\&]
\Pi_1P^*_{d_0\deleted{\unsure{\geq}}d_1,\ldots,d_r}\arrow[r]\&
\Pi_1P^*_{d'_0,\ldots,d'_r}
\end{tikzcd}}
\margin{Attention, on a subrepticement changé de notations, en intervertissant $d_*$ et $d'_*$ !}
\margin{on se réduit au cas de $d_0,d_1,\ldots,d_r$ strict.t décroissants -- auquel on parvient aussitôt. En fait, on est ramené au cas où $\unsure{r'=r-1}$, et où il s'agit de l'opération « face de codim $1$ » pour laquelle $\alpha(0)\ne0$, donc correspondant :
\diagram{\begin{tikzcd}[ampersand replacement=\&]D_{d_0,d_1,\ldots,d_r}\arrow[r]\&D_{d_1,\ldots,d_r}\end{tikzcd}}}
Notons que l'on a
\[
(223)
\]
\diagram{\begin{tikzcd}[ampersand replacement=\&]
D^{\deleted{*}}_{d_0,\ldots,d_r}\arrow[r]\&D_{d'_0,\ldots,d'_{r'}}
\end{tikzcd}}
(mais \uline{non} un morphisme entre les
espaces avec $*$ !), qui est une immersion \inserted{locale}
d'espaces analytiques lisses, de codim
$d_0-d'_0$, \deleted{qui}\inserted{\unsure{ainsi|et}} s'insère dans le diagramme
\margin{134}
